% -----------------------------------------------------------------------------
% -----------------------------------------------------------------------------

\begin{exbox}[四点 Green 函数到截面]
取稳定的实标量粒子，使用树级留数 $Z=1$，相互作用为
\begin{equation*}
 \Lag_{\rm int}=-\frac{g_4}{4!\,\hbar c}\varphi^4,
 \qquad [g_4]=1,
\end{equation*}
物理外态的四动量满足 $p_i^2=m^2c^2$。从四点连通函数出发，用 LSZ 约化求在壳振幅，再计算二体散射的角分布与总截面。

\textbf{解。} 在动量空间中，四点连通函数的最低阶奇异部分为
\begin{equation}
 \begin{aligned}
 \widetilde G_{4,c}^{(1)}
 =&\,(2\pi\hbar)^4\delta^{(4)}(p_1+p_2-p_3-p_4)
 \left(-\frac{\ii g_4}{\hbar^2c^2}\right)\\
 &\times\prod_{a=1}^{4}
 \frac{\ii\hbar^3c}{p_a^2-m^2c^2+\ii0^+}.
 \end{aligned}
 \label{eq:phi4-fourpoint-poles-dp64}
\end{equation}
四维 $\delta$ 函数来自顶角位置积分，中间常数来自局域四价顶角，四个简单极点来自相互作用点到外部场插入点的传播。LSZ 约化把这些场插入转换成渐近单粒子外态，因此需要消去外腿极点。

对第 $a$ 条外腿定义逆极点因子
\begin{equation*}
 R_a(p_a)=\frac{p_a^2-m^2c^2}{\ii\hbar^3c}.
\end{equation*}
它与相应传播子在在壳留数极限下相消：
\[
 \lim_{p_a^2\to m^2c^2}
 R_a(p_a)\frac{\ii\hbar^3c}{p_a^2-m^2c^2+\ii0^+}=1.
\]
对四条外腿分别实施这一约化后，保留的顶角与总动量守恒因子为
\begin{equation*}
 (2\pi\hbar)^4\delta^{(4)}(P_f-P_i)
 \left(-\frac{\ii g_4}{\hbar^2c^2}\right).
\end{equation*}
按前文的国际单位制归一化约定，提出公共场归一化因子 $1/(\hbar^2c^2)$，余下系数满足 $\ii\mathcal M=-\ii g_4$，所以在壳约化振幅为
\begin{equation}
 \mathcal M_{\varphi\varphi\to\varphi\varphi}=-g_4.
 \label{eq:phi4-reduced-amplitude-dp64}
\end{equation}
外腿传播子已从振幅中消去，四个外粒子则由在壳动量标签标识。

将\eqnrefs{eq:phi4-reduced-amplitude-dp64} 代入\eqnrefs{eq:2to2-cross-section-master-si-r9}。等质量弹性散射满足 $|\vb{p}_f|=|\vb{p}_i|$，末态两个标量完全相同，因此全立体角计数需要除以 $2!$，得到
\begin{equation*}
 \frac{\dd\sigma}{\dd\Omega}
 =\frac12\frac{\hbar^2}{64\pi^2s}g_4^2
 =\frac{\hbar^2g_4^2}{128\pi^2s},
\end{equation*}
接触顶角没有 $t$ 或 $u$ 依赖，所以角分布各向同性。用 $\int\dd\Omega=4\pi$ 积分整个立体角，得到
\begin{equation}
 \sigma_{\rm tot}
 =\frac{\hbar^2g_4^2}{32\pi s}.
 \label{eq:phi4-total-cross-dp64}
\end{equation}
若用质心能量 $E_{\rm CM}=c\sqrt{s}$ 表示，便有
\begin{equation*}
 \sigma_{\rm tot}
 =\frac{g_4^2}{32\pi}
 \left(\frac{\hbar c}{E_{\rm CM}}\right)^2.
\end{equation*}
这里 $\hbar c/E_{\rm CM}$ 具有长度量纲，其平方给出截面的量纲。在这一树级接触模型中，固定 $g_4$ 时截面按 $E_{\rm CM}^{-2}$ 衰减。作为数量级估算，取 $g_4=0.50$、$E_{\rm CM}=1.00\,\mathrm{GeV}$，用 $\hbar c=0.19732698\,\mathrm{GeV\,fm}$ 得
\begin{align*}
 \sigma_{\rm tot}
 &=\frac{0.50^2}{32\pi}
 \left(\frac{0.19732698\,\mathrm{GeV\,fm}}{1.00\,\mathrm{GeV}}\right)^2,\\
 &\simeq9.68\times10^{-5}\,\mathrm{fm^2}
 =9.68\times10^{-35}\,\mathrm{m^2}
 \simeq0.968\,\mu\mathrm b.
\end{align*}
此数值来自同一振幅与相空间，不增加模型假设。全同末态因子、国际单位制归一化和树级截断都保留在计算中；加入圈修正后，振幅及其能量依赖还会改变。
\end{exbox}




